\documentclass[12pt]{article}
\usepackage[T1]{fontenc}
\usepackage[utf8]{inputenc}
\usepackage{lmodern}
\usepackage{textcomp}
\usepackage{enumitem}
\usepackage{geometry}
\geometry{margin=1in}
\begin{document}
\begin{center}
Answer Key - Sociology - Graduate Level Assessment 7 \
\small (Answers are ordered exactly as in the question paper.)
\end{center}

\section*{Multiple Choice}
\begin{enumerate}[label=\arabic*.]
\item \textbf{Answer:} C
\\ \textit{Options:} A) commuters started moving out of villages and into cities; B) towns and cities were becoming increasingly planned and managed; C) industrial capitalism led to a shift of population from rural to urban areas; D) transport systems were not provided, so it was easier to live in the city
\\ \textit{Note:} The correct answer highlights that industrial capitalism created job opportunities in cities, prompting a significant migration of people from rural areas seeking employment and better living conditions during the urbanization of the nineteenth century.
\item \textbf{Answer:} C
\\ \textit{Options:} A) the content of the media is determined by market forces; B) the subordinate classes are dominated by the ideology of the ruling class; C) the media manipulate 'the masses' as vulnerable, passive consumers; D) audiences make selective interpretations of media messages
\\ \textit{Note:} Mass-society theory posits that the media play a crucial role in shaping and controlling the beliefs and behaviors of the masses, portraying them as passive recipients who are easily influenced by media messages, thus highlighting the manipulative power of media in a society where individuals feel isolated and disconnected.
\item \textbf{Answer:} B
\\ \textit{Options:} A) contrasted homosexuality with heterosexuality; B) distinguished between male and females as separate sexes; C) represented women's genitalia as underdeveloped versions of men's; D) argued for male superiority over women
\\ \textit{Note:} The 'two-sex' model as identified by Laqueur emphasizes the historical and sociocultural distinction between male and female as separate entities, challenging previous notions of a more fluid, singular understanding of sex.
\item \textbf{Answer:} A
\\ \textit{Options:} A) highly specialized, interrelated sets of social practices; B) disorganized social relations in a postmodern world; C) virtual communities in cyberspace; D) no longer relevant to sociology
\\ \textit{Note:} The correct answer highlights that in contemporary societies, social institutions have evolved into highly specialized entities that are interconnected through various social practices, reflecting the complexity and interdependence of societal functions.
\item \textbf{Answer:} A
\\ \textit{Options:} A) ways of acting, thinking, and feeling that are collective and social in origin; B) the way scientists construct knowledge in a social context; C) data collected about social phenomena that are proven to be correct; D) ideas and theories that have no basis in the external, physical world
\\ \textit{Note:} This answer is correct because Durkheim emphasized that social facts are phenomena that exist outside of the individual and are shaped by the collective norms, values, and beliefs of society, thereby influencing individual behavior and thought.
\end{enumerate}

\section*{True / False}
\begin{enumerate}[label=\arabic*.]
\setcounter{enumi}{5}
\item \textbf{Answer:} False
\\ \textit{Options:} A) True; B) False
\\ \textit{Note:} In the 19 th century, employers primarily exercised disciplinary power through stringent oversight, harsh working conditions, and punitive measures rather than by making routine tasks less monotonous for workers.
\item \textbf{Answer:} False
\\ \textit{Options:} A) True; B) False
\\ \textit{Note:} The statement is false because the post-war welfare state of 1948 primarily focused on social security and economic stability rather than fully ensuring free health care and education for all citizens, which developed more comprehensively in subsequent decades.
\item \textbf{Answer:} False
\\ \textit{Options:} A) True; B) False
\\ \textit{Note:} The statement is false because Marx's concept of 'mode of production' encompasses not only the methods of production in factories but also the social relations, economic systems, and historical contexts that shape and are shaped by the production processes.
\item \textbf{Answer:} True
\\ \textit{Options:} A) True; B) False
\\ \textit{Note:} Sociology is considered a social science because it employs systematic methodologies to develop theories that are grounded in observable data and rigorous analysis, ensuring that its conclusions are both logical and empirically validated.
\item \textbf{Answer:} False
\\ \textit{Options:} A) True; B) False
\\ \textit{Note:} Marx did not suggest a specific percentage threshold for church attendance to determine the cessation of religion; rather, he viewed religion as a social construct tied to material conditions and class struggles, which would evolve rather than simply disappear at a certain attendance level.
\end{enumerate}

\section*{Long Form Response}
\begin{enumerate}[label=\arabic*.]
\setcounter{enumi}{10}
\item \textbf{Answer:} Wallerstein's concept of the 'world system' situates the capitalist world economy as a complex, interconnected system that transcends national boundaries, fostering a global division of labor that actively includes diverse ethnic and cultural groups. This framework highlights how labor is organized not only by economic factors but also by historical and social contexts, leading to distinct roles for various groups within the global economy. For instance, certain ethnic groups may be relegated to low-wage, labor- intensive jobs, while others dominate higher-value sectors, illustrating the persistence of inequality and exploitation. Consequently, Wallerstein's analysis underscores the implications of capitalism in shaping social hierarchies and cultural identities, as the world system perpetuates divisions that are both economic and cultural in nature, influencing the lived experiences of individuals across the globe.
\\ \textit{Note:} correct because it accurately describes Wallerstein's 'world system' theory as a framework that illustrates how the capitalist world economy creates a global division of labor influenced by historical, social, and economic factors, thereby assigning specific roles to diverse ethnic and cultural groups.
\item \textbf{Answer:} The educational policies of the New Labour government, particularly in the late 1990 s and early 2000 s, emphasized the need for reform and improvement in the education sector, often at the expense of failing Local Education Authorities (LEAs). By not supporting these LEAs in their reform efforts, the government risked perpetuating educational inequalities, as schools in disadvantaged areas frequently lacked the resources and support necessary for effective change. This approach also led to a reliance on market mechanisms, which favored well-performing schools and created a competitive environment that marginalized those in need of significant support. Consequently, the decision to withhold assistance from failing LEAs not only undermined potential educational reforms but also reinforced systemic disparities, ultimately affecting student outcomes and community cohesion. Thus, the implications of this policy illuminate the broader sociological impact of neglecting local governance in education, highlighting the need for a more equitable and supportive framework for all LEAs.
\\ \textit{Note:} correct because it highlights how the New Labour government's lack of support for failing LEAs contributed to educational inequalities by leaving disadvantaged schools without the necessary resources for effective reform.
\end{enumerate}

\vfill
\noindent\textit{End of Answer Key}
\end{document}